% This chapter is the complete working manuscript for the thirteenth seminar
% talk. The final section contains supplementary source material outside the
% default live route.

\chapter{Elliptic representability II: derived charts and globalization}
\label[chapter]{chap:Elliptic_Representability_Derived_Charts}

The previous chapter ended with one precise statement. Near a chosen solution
of an elliptic differential moduli problem, the augmented operator
\[
  \widetilde P\colon
  \widetilde U\longrightarrow S\times\Gamma(E;N)
\]
is submersive as a morphism of derived $\Cinfty$-stacks. All
infinite-dimensional analysis was concentrated in the proof of that
statement.

This chapter turns it into geometry. We first take the zero fiber of
$\widetilde P$. Submersiveness makes this an ordinary finite-dimensional
manifold $V$. We then project $V$ to the finite-dimensional obstruction space
and take one further derived zero fiber. Pullback pasting identifies the
result with an open restriction of the intrinsic solution stack. Varying the
solution produces an effective open atlas and completes the proof of
\Cref{thm:Relative_Elliptic_Representability}.

The final part compares this mechanism with Pardon's organization of the same
Fredholm geometry. Steffens proves derived submersiveness of an augmented
operator through a Sobolev tower. Pardon first represents the ordinary
regular locus, extends proper section stacks from smooth to derived test
objects, and makes an arbitrary solution regular by enlarging the parameter
space. The comparison is structural. Pardon's published article gives a
proof sketch for pseudo-holomorphic section problems, not a second complete
proof of Steffens's theorem in its full generality.

\section{The finite-dimensional derived chart}
\label[section]{sec:Finite_Dimensional_Derived_Chart_Thirteen}

\subsection{The augmented zero fiber is ordinary}
\label[subsection]{sec:Augmented_Zero_Fiber_Ordinary_Thirteen}

Retain the local data from the previous chapter. Thus $S$ is a manifold,
$N$ is compact,
\[
  P\colon S\times Q\longrightarrow S\times\Gamma(E;N)
\]
is the local family of elliptic equations, and
\[
  \widetilde U\subseteq S\times Q\times C
\]
is an augmented neighborhood on which
\begin{equation}
  \label{eq:Augmented_Operator_Thirteen}
  \widetilde P(s',\tau,c)
  =\bigl(s',P_{s'}(\tau)+\iota(c)\bigr)
\end{equation}
is derived-submersive.

Let $0_E\colon S\to S\times\Gamma(E;N)$ be the zero section and form the
derived pullback
\begin{equation}
  \label{eq:Augmented_Zero_Fiber_Thirteen}
  \begin{tikzcd}
    V \rar \dar \drar[pullback]
      & \widetilde U \dar["\widetilde P"] \\
    S \rar["0_E"']
      & S\times\Gamma(E;N).
  \end{tikzcd}
\end{equation}
By \Cref{prop:Augmented_Operator_Derived_Submersive_Twelve}, the right-hand
map is submersive. The square is therefore transverse, so its derived
pullback is represented by the ordinary smooth pullback
\begin{equation}
  \label{eq:Augmented_Solution_Manifold_Thirteen}
  V
  =\{(s',\tau,c)\in\widetilde U:
       P_{s'}(\tau)+\iota(c)=0\}.
\end{equation}
This is a finite-dimensional manifold, and $V\to S$ is a submersion.

The first pullback has removed the infinite-dimensional ambient section
spaces. It has not yet imposed the original equation, because a point of $V$
may have $c\neq0$.

\subsection{The obstruction section}
\label[subsection]{sec:Relative_Obstruction_Section_Thirteen}

Projection to the obstruction coordinate restricts to a map over $S$,
\begin{equation}
  \label{eq:Obstruction_Section_Thirteen}
  \kappa\colon V\longrightarrow S\times C,
  \qquad
  (s',\tau,c)\longmapsto(s',c).
\end{equation}
Let $0_C\colon S\to S\times C$ be the zero section. Define
\begin{equation}
  \label{eq:Derived_Obstruction_Zero_Locus_Thirteen}
  \begin{tikzcd}
    \mathcal Z_t \rar \dar \drar[pullback]
      & V \dar["\kappa"] \\
    S \rar["0_C"']
      & S\times C.
  \end{tikzcd}
\end{equation}
On ordinary points, $c=0$ reduces
$P_{s'}(\tau)+\iota(c)=0$ to $P_{s'}(\tau)=0$. The derived statement is
stronger and follows by pasting pullback squares.

To make the pasting precise, let
\[
  U_0
  =\widetilde U\times_{S\times Q\times C}(S\times Q),
\]
where $S\times Q\to S\times Q\times C$ is the zero section. The map
$U_0\to S\times Q$ is an open inclusion and
$\widetilde P|_{U_0}=P|_{U_0}$. Pasting
\Cref{eq:Augmented_Zero_Fiber_Thirteen} and
\Cref{eq:Derived_Obstruction_Zero_Locus_Thirteen} gives an isomorphism
\begin{equation}
  \label{eq:Derived_Chart_Local_Solution_Isomorphism_Thirteen}
  \mathcal Z_t
  \iso
  U_0\times_{S\times\Gamma(E;N)}S.
\end{equation}
The right side is the derived zero locus of $P$ restricted to $U_0$. By
\Cref{eq:Local_Solution_Open_Substack_Twelve} and
\Cref{lem:Open_Charts_Pass_To_Pullbacks_Twelve}, it is an open restriction
of the intrinsic stack $\Sol(\mathcal E)$ containing $t$.

This argument is the decisive comparison. Equality of ordinary zero sets and
agreement of tangent complexes would not identify the derived moduli
functors. The pasted universal properties do.

\subsection{Quasi-smoothness and the tangent complex}
\label[subsection]{sec:Quasi_Smoothness_Tangent_Thirteen}

Both $V$ and $S\times C$ are finite-dimensional manifolds over $S$. After
shrinking around $t$, they are affine $\Cinfty$-schemes of finite
presentation. The derived zero locus in
\Cref{eq:Derived_Obstruction_Zero_Locus_Thirteen} is therefore a quasi-smooth
derived $\Cinfty$-scheme locally of finite presentation.

At $v=(s,\sigma,0)\in V$, its relative tangent complex is
\begin{equation}
  \label{eq:Kuranishi_Tangent_Complex_Thirteen}
  \left[
    T_v(V/S)\xrightarrow{D_v\kappa}C
  \right]
\end{equation}
in cohomological degrees $0$ and $1$. We compare it directly with the
elliptic complex. Put
\[
  L=D_\sigma P_s\colon
  \Gamma(F;N)\longrightarrow\Gamma(E;N).
\]
The derivative of the augmented equation is
\[
  A=(L,\iota)\colon
  \Gamma(F;N)\oplus C
  \longrightarrow
  \Gamma(E;N),
\]
and it is surjective. Moreover,
\[
  T_v(V/S)=\ker A,
  \qquad
  D_v\kappa=\operatorname{pr}_C|_{\ker A}.
\]
There is therefore a chain map
\begin{equation}
  \label{eq:Kuranishi_Elliptic_Chain_Map_Thirteen}
  \begin{tikzcd}[column sep=large]
    \ker A \rar["\operatorname{pr}_C"] \dar["\operatorname{pr}_{\Gamma(F)}"']
      & C \dar["-\iota"] \\
    \Gamma(F;N) \rar["L"']
      & \Gamma(E;N).
  \end{tikzcd}
\end{equation}
It commutes because $L(\xi)+\iota(c)=0$ for $(\xi,c)\in\ker A$.

\begin{proposition}[Finite and elliptic tangent complexes agree]
  \label[proposition]{prop:Finite_Elliptic_Tangent_Complexes_Agree}
  The chain map in
  \Cref{eq:Kuranishi_Elliptic_Chain_Map_Thirteen} is a quasi-isomorphism.
  Consequently,
  \[
    \mathbb T_{\mathcal Z_t/S,v}
    \simeq
    [\Gamma(F;N)\xrightarrow{L}\Gamma(E;N)].
  \]
\end{proposition}

\begin{proof}
  On degree-zero cohomology, the map identifies
  \[
    \ker(D_v\kappa)
    =\{(\xi,0):L\xi=0\}
  \]
  with $\ker L$.

  For degree one, the induced map is
  \[
    C/\operatorname{im}(D_v\kappa)
    \longrightarrow
    \operatorname{coker}L,
    \qquad
    [c]\longmapsto[-\iota(c)].
  \]
  It is surjective because $A=(L,\iota)$ is surjective. Its kernel consists
  of those $c$ for which $\iota(c)=-L\xi$ for some $\xi$. Equivalently,
  $(\xi,c)\in\ker A$, so $c$ lies in the image of $D_v\kappa$. Hence the map
  is an isomorphism.
\end{proof}

The two pullbacks have complementary roles:
\[
  \begin{gathered}
    \widetilde P\text{ is submersive}
    \quad\Longrightarrow\quad
    V\text{ is an ordinary manifold},\\
    \kappa\text{ need not be submersive}
    \quad\Longrightarrow\quad
    \mathcal Z_t\text{ may carry derived structure}.
  \end{gathered}
\]
The augmentation removes the infinite-dimensional failure of transversality.
The derived zero locus retains the remaining finite-dimensional obstruction.

\section{Globalization and the representability theorem}
\label[section]{sec:Globalization_Representability_Thirteen}

\subsection{An effective open atlas}
\label[subsection]{sec:Effective_Open_Atlas_Thirteen}

The construction above applies near every solution $t$. For each such point,
we obtain a quasi-smooth derived $\Cinfty$-scheme $\mathcal Z_t$ locally of
finite presentation and an open morphism
\[
  \mathcal Z_t\longrightarrow\Sol(\mathcal E)
\]
whose image contains $t$.

These maps cover the solution stack. More precisely, every family of
solutions over a test manifold is locally on that manifold contained in one
of the chosen section-stack charts and augmented neighborhoods. Thus
\begin{equation}
  \label{eq:Effective_Open_Atlas_Thirteen}
  \coprod_t\mathcal Z_t
  \longrightarrow
  \Sol(\mathcal E)
\end{equation}
is an effective epimorphism and each summand is an open substack.

There is no independent coordinate-change problem to solve. On an overlap,
both charts represent open restrictions of the same intrinsic stack. Their
intersection is the corresponding pullback over $\Sol(\mathcal E)$, and its
coherent transition data are inherited from that stack.

\subsection{Completion of the proof}
\label[subsection]{sec:Completion_Relative_Representability_Thirteen}

\begin{proof}[Proof of \Cref{thm:Relative_Elliptic_Representability}]
  Work first after a representable base change to a manifold chart of $S$.
  Around every solution, the localization in Chapter~12 gives one smooth
  family of elliptic equations over a compact manifold. The Sobolev-scale
  argument and
  \Cref{prop:Augmented_Operator_Derived_Submersive_Twelve} make its augmented
  operator derived-submersive.

  The two pullbacks in
  \Cref{eq:Augmented_Zero_Fiber_Thirteen,eq:Derived_Obstruction_Zero_Locus_Thirteen}
  produce a quasi-smooth derived $\Cinfty$-scheme locally of finite
  presentation. Pullback pasting identifies it with an open substack of the
  intrinsic solution stack. Varying the solution yields the effective open
  atlas in \Cref{eq:Effective_Open_Atlas_Thirteen}.

  Apply the open-atlas criterion
  \Cref{thm:Open_Local_Representability}. It follows that the solution stack
  after base change to the manifold chart is represented. Finally, use the
  relative consequence
  \Cref{cor:Local_Representability_Relative} to descend this conclusion over
  the original smooth-stack base $S$. Quasi-smoothness and local finite
  presentation are local on the representing scheme, so they descend as
  well.
\end{proof}

This final globalization is deliberately short. The technical work was the
construction and identification of one open chart. Once those charts are
open substacks of a single intrinsic stack, Proposition~2.1.48 and
Corollary~2.1.49 perform the gluing
\cite[Proposition~2.1.48 and Corollary~2.1.49]{Steffens_Representability_Elliptic_Moduli}.

\subsection{What the hypotheses and choices do}
\label[subsection]{sec:Hypotheses_Choices_Thirteen}

The proof assigns a distinct role to every hypothesis.
\begin{enumerate}
  \item The smooth-stack base permits descent to manifold charts while
    retaining families and automorphisms.
  \item Properness gives compact fibers and local triviality of the family of
    domains.
  \item Finite order makes the nonlinear equation functorial through relative
    jets and compatible with the Sobolev tower.
  \item Ellipticity gives Fredholm linearizations and the regularity which
    returns Sobolev solutions to smooth sections.
  \item Finite-dimensional obstruction spaces turn the local problem into a
    quasi-smooth finite-presentation derived zero locus.
\end{enumerate}

Local trivializations, local additions, Sobolev indices, complements, and
representatives of the cokernel all occur in the construction. They do not
occur in the definition of $\Sol(\mathcal E)$. Every resulting chart is
identified with an open restriction of that same intrinsic stack. This is the
precise sense in which the representing object is independent of the
analytic choices.

\section{The proof pattern in finite dimensions}
\label[section]{sec:Proof_Pattern_Finite_Dimensions_Thirteen}

\subsection{The complete route}
\label[subsection]{sec:Complete_Representability_Route_Thirteen}

The proof can be read as the following chain of constructions:
\[
  \begin{aligned}
    &\text{Intrinsic elliptic solution stack}\\
    &\quad\rightsquigarrow
      \text{local smooth nonlinear operator }P\\
    &\quad\rightsquigarrow
      \text{derived-submersive augmentation }\widetilde P\\
    &\quad\rightsquigarrow
      \text{ordinary manifold }V\xrightarrow{\kappa}S\times C\\
    &\quad\rightsquigarrow
      V\times_{S\times C}^{\mathbf R}S\\
    &\quad\rightsquigarrow
      \text{effective open atlas}.
  \end{aligned}
\]
The first two arrows use differential geometry and jets. The third uses
Fredholm analysis, the Sobolev tower, and the smooth-to-derived criterion.
The fourth is a transverse pullback. The fifth is finite-dimensional derived
geometry. The last uses descent.

\subsection{The model calculation}
\label[subsection]{sec:X_Squared_Regular_Enlargement_Thirteen}

The whole route is visible in the function
\[
  f\colon\R\longrightarrow\R,
  \qquad
  f(x)=x^2.
\]
Its ordinary zero locus is one point, but its derived zero locus has tangent
complex
\[
  [\R\xrightarrow{0}\R]
\]
at the origin. Add an obstruction coordinate and define
\[
  \widetilde f\colon\R^2\longrightarrow\R,
  \qquad
  \widetilde f(x,c)=x^2+c.
\]
This map is a submersion. Its zero fiber
\[
  V=\{(x,c):x^2+c=0\}
\]
is an ordinary one-dimensional manifold. Restrict the second coordinate to
obtain
\[
  \kappa\colon V\longrightarrow\R,
  \qquad
  \kappa(x,-x^2)=-x^2.
\]
Pullback pasting gives
\begin{equation}
  \label{eq:X_Squared_Regular_Enlargement_Thirteen}
  \R\times_{\R}^{\mathbf R}\{0\}
  \iso
  V\times_{\R}^{\mathbf R}\{0\},
\end{equation}
where the map on the left is $f$ and the map on the right is $\kappa$.

The left side presents the singular equation directly. The right side first
solves a regular enlarged equation and then imposes the obstruction parameter
by a derived fiber. Steffens's proof is the elliptic family version of this
calculation.

\subsection{Representability is not a virtual class}
\label[subsection]{sec:Representability_Not_Virtual_Class_Thirteen}

The theorem proves that the moduli functor is represented and that its local
derived geometry is quasi-smooth. It does not automatically produce an
enumerative invariant. A virtual fundamental class additionally requires a
suitable global setting, typically including compactness or properness of the
moduli object and orientation data for its determinant complex.

The bordism interlude explained what a compact oriented derived manifold can
contribute. The present theorem supplies a represented quasi-smooth object,
but it does not establish those further hypotheses. This distinction remains
essential in applications.

\section{Comparison with Pardon's regular-locus engine}
\label[section]{sec:Pardon_Regular_Locus_Comparison_Thirteen}

\subsection{The regular locus on smooth tests}
\label[subsection]{sec:Pardon_Regular_Locus_Smooth_Tests_Thirteen}

Pardon's published argument begins with a pseudo-holomorphic section problem
over a proper family of smooth compact curves. Schematically, one has an
equation morphism
\[
  \Phi\colon
  \operatorname{Sect}_{C/B}(W)
  \longrightarrow
  \operatorname{Sect}_{C/B}(H)
\]
and the solution stack is its fiber over the zero section.

The \emph{regular locus} consists of the solutions where the total
linearization of $\Phi$ is surjective. Ordinary parametric Fredholm theory
represents this locus, on smooth test manifolds, by an ordinary smooth
manifold. The new question is not analytic:
\begin{quote}
  Why does the same manifold classify regular families over derived test
  objects?
\end{quote}

\subsection{Proper section stacks and derived tests}
\label[subsection]{sec:Pardon_Proper_Section_Stacks_Thirteen}

Let
\[
  i\colon\mathrm{Sm}\longrightarrow\mathrm{Der}
\]
denote the inclusion of manifolds into derived smooth manifolds. Restriction
of stacks has a left adjoint $i_!$, the sheaf left Kan extension. For a stack
$F$ on manifolds, $i_!F$ is the universal extension of $F$ to derived tests
which respects the required colimits and descent.

Pardon's Proposition~5.2 says, in the form needed here, that the section stack
of a submersion over a proper smooth source is obtained by this universal
extension:
\begin{equation}
  \label{eq:Pardon_Proper_Section_Stack_Thirteen}
  i_!\operatorname{Sect}_{C/B}(W)_{\mathrm{Sm}}
  \iso
  \operatorname{Sect}_{C/B}(W)_{\mathrm{Der}}.
\end{equation}
Properness and the softness supplied by partitions of unity are the
geometric inputs. The statement does not say that derived families are
ordinary families. It says that their coherent behavior on derived tests is
determined by the ordinary smooth functor.

Left adjoints do not preserve arbitrary pullbacks. Pardon's Lemma~5.3 gives
the additional fact used in the proof: sheaf left Kan extension preserves the
pullbacks in which the distinguished morphism is submersive. On the regular
locus, the equation morphism is submersive. Applying the lemma to its zero
fiber shows that
\begin{equation}
  \label{eq:Pardon_Regular_Locus_Extension_Thirteen}
  i_!\operatorname{Sol}(\Phi)^{\mathrm{reg}}_{\mathrm{Sm}}
  \iso
  \operatorname{Sol}(\Phi)^{\mathrm{reg}}_{\mathrm{Der}}.
\end{equation}
Thus the ordinary manifold which represents the regular locus on smooth tests
also represents it on derived tests
\cite[Proposition~5.2 and Lemma~5.3]{Pardon_Representability_Elliptic_Fredholm}.

\subsection{Parameter thickening}
\label[subsection]{sec:Pardon_Parameter_Thickening_Thirteen}

At a nonregular solution, choose a finite-dimensional space $C$ mapping
surjectively to the cokernel of the linearization. Enlarge the parameter base
by $C$ and add the corresponding inhomogeneous terms to the equation. The
chosen solution is regular in this enlarged family. The preceding argument
represents the enlarged regular locus by an ordinary manifold. The original
problem is recovered by taking its derived fiber over the zero value of the
new parameter.

This is the same enlargement-and-derived-fiber pattern as
\Cref{eq:X_Squared_Regular_Enlargement_Thirteen}. The two proofs organize its
middle step differently:
\begin{enumerate}
  \item Steffens adds an obstruction coordinate to the operator and uses the
    functorial Sobolev tower to prove that the augmented smooth operator is
    derived-submersive.
  \item Pardon enlarges the parameter base, represents the ordinary regular
    locus, uses proper section stacks to extend it to derived tests, and
    returns by a derived parameter fiber.
\end{enumerate}

Both methods obtain local representatives as open substacks of an intrinsic
solution functor. Consequently, both avoid constructing a strict cocycle of
independently chosen Kuranishi charts.

\subsection{Scope and source status}
\label[subsection]{sec:Pardon_Scope_Source_Status_Thirteen}

Pardon's published Theorem~5.1 concerns pseudo-holomorphic section problems
over proper families of smooth compact curves. Its conclusion and hypotheses
are not identical to Steffens's general theorem on elliptic differential
moduli problems. The article explicitly presents the argument as a brief
proof sketch
\cite[Theorem~5.1 and its proof sketch]{Pardon_Representability_Elliptic_Fredholm}.

Pardon's July 2026 manuscript develops a much broader logarithmic and
elliptic framework and gives a longer account of the regular-locus strategy.
It also explicitly identifies itself as unfinished work in progress and
contains unresolved internal placeholders. In particular, the proposed
logarithmic derived-base reduction is not a completed proof source. We use
the manuscript only to clarify the architecture of the published sketch
\cite[Section~N.6.4]{Pardon_Logarithmic_Derived_Moduli}.

\section{Summary and supplementary materials}
\label[section]{sec:Derived_Charts_Globalization_Summary_Supplementary}

\subsection{Takeaways}
\label[subsection]{sec:Derived_Charts_Globalization_Takeaways}

The live route has nine conclusions.
\begin{enumerate}
  \item The zero fiber of the derived-submersive augmented operator is an
    ordinary finite-dimensional manifold $V$.
  \item Projection to the obstruction coordinate gives a finite-dimensional
    map $\kappa\colon V\to S\times C$.
  \item Its derived zero locus recovers the original local equation.
  \item Pullback pasting identifies this derived zero locus with an open
    restriction of the intrinsic solution stack.
  \item Its tangent complex is quasi-isomorphic to the elliptic
    deformation-obstruction complex.
  \item The local chart is quasi-smooth and locally of finite presentation.
  \item The local charts form an effective open atlas, so the open-atlas
    theorem proves relative representability.
  \item The analytic choices affect presentations, not the intrinsic
    represented stack.
  \item Pardon's regular-locus proof uses the same enlargement-and-derived-
    fiber pattern with a different smooth-to-derived engine.
\end{enumerate}

\subsection{Supplement: the precise published Pardon endpoint}
\label[subsection]{sec:Supplement_Pardon_Published_Endpoint_Thirteen}

Pardon's published Derived Regularity Theorem states representability for a
pseudo-holomorphic section problem
\[
  W\longrightarrow C\longrightarrow B
\]
over a derived smooth stack, with $C\to B$ a proper family of real
two-dimensional domains and with the relevant complex structures. It also
identifies the derived-to-topological comparison for the solution stack
\cite[Theorem~5.1]{Pardon_Representability_Elliptic_Fredholm}.

The printed statement says that $B\to C$ is proper of relative dimension two,
although the displayed composable maps and the proof require $C\to B$. This
is best treated as a typographical reversal, not as a new mathematical
hypothesis.

\subsection{Related literature}
\label[subsection]{sec:Derived_Charts_Globalization_Related_Literature}

Steffens's Theorem~3.4.3 is the primary source for the completed
representability proof. The local chart and globalization occupy the final
part of its proof, while Proposition~2.1.48 and Corollary~2.1.49 supply the
open-atlas step
\cite[Theorem~3.4.3, Proposition~2.1.48, and
Corollary~2.1.49]{Steffens_Representability_Elliptic_Moduli}.

Pardon's published proceedings article supplies the comparison through
Theorem~5.1, Proposition~5.2, and Lemma~5.3
\cite[Section~5]{Pardon_Representability_Elliptic_Fredholm}. The July 2026
manuscript is useful for terminology and proof order, but its work-in-progress
status must be retained whenever it is cited
\cite[Sections~P.5 and N.6.4]{Pardon_Logarithmic_Derived_Moduli}.
